\chapter{Factor Endowments}\label{ch:b3:factor-endowments}

Ricardo's model explains trade by differences in technology. Countries also
differ in what they have: some have much capital per worker, some many workers
with little schooling, some vast farmland. When the same technology is available
everywhere, these differences in endowments are enough to create \omterm{def:b3:ricardian-trade:comparative-advantage}{comparative
advantage}. This chapter builds the models in which they do --- first with
factors stuck in their sectors, then with factors that move --- and draws from
them the two predictions economists use most: who gains and who loses from
trade inside a country, and which factors a country's trade embodies. The second
prediction has been tested for seventy years, and the record of those tests is
the chapter's last section.

\begin{openingchart}{Shares of the world's capital stock, workers and output in
2019, for the twelve largest economies by output (percent; the world is the 175
countries with all three series).}{Penn World Table 11.0; our computation}
\begin{tikzpicture}
\begin{axis}[omscatter, width=0.72\linewidth, height=0.55\linewidth, xmin=0, xmax=25,
  ymin=0.4, ymax=12.6, ytick={1,...,12}, yticklabels from table={figdata/bachelor-3/08-factor-endowments/out/endow.csv}{country},
  y tick label style={font=\footnotesize}, xlabel={Share of the world total (\%)},
  legend pos=south east, ymajorgrids=true, xmajorgrids=false]
  \addplot[only marks, mark=square*, mark size=2.2pt, omDef] table[x=capital, y=rank] {figdata/bachelor-3/08-factor-endowments/out/endow.csv};
  \addplot[only marks, mark=triangle*, mark size=2.6pt, omThm] table[x=workers, y=rank] {figdata/bachelor-3/08-factor-endowments/out/endow.csv};
  \addplot[only marks, mark=*, mark size=2pt, black] table[x=output, y=rank] {figdata/bachelor-3/08-factor-endowments/out/endow.csv};
  \legend{capital, workers, output}
\end{axis}
\end{tikzpicture}
\end{openingchart}

The United States had 11.4~\% of the world's capital and 4.9~\% of its workers;
India 6.1~\% of the capital and 16.2~\% of the workers; China, the largest in
both, had more of the workers (23.0~\%) than of the capital (19.1~\%).
% figdata: endow
Japan, Germany, France, the United Kingdom, Italy and South Korea all hold two to
four times as large a share of the world's capital as of its workers. If
technology were the same everywhere, these ratios alone would tell each country
what to export.

\begin{recall}
The grades 10--12 volume drew the production possibility frontier and its
opportunity costs. The Year 1 volume taught production functions, isoquants and
cost minimisation, and the Year 2 volume general equilibrium. Here they are
combined, two goods and two factors at a time.
\end{recall}

\begin{notation}[Symbols of this chapter]\label{not:b3:factor-endowments:symbols}
Goods $j = 1, 2$ have prices $p_j$ (not the price level). Labour earns the wage
$w$ and capital the rental rate $r$ (not an interest rate). $a_{Lj}$ and $a_{Kj}$
are the labour and capital used per unit of good $j$, and $\theta_{Lj} = w
a_{Lj}/p_j$ is labour's share of its cost. A hat is a proportional change:
$\hat x = dx/x$. In the factor-content section, $A$ is a matrix of factor
requirements (not productivity), $V_c$ a country's endowments, $T_c$ its net
exports and $s_c$ its share of world spending.
\end{notation}

\section{Specific factors: who gains in the short run}\label{sec:b3:factor-endowments:specific}

\begin{definition}[Specific-factors model]\label{def:b3:factor-endowments:specific-factors}
\index{specific-factors model}
In the \emph{specific-factors model}, two goods are made with labour, which moves
freely between sectors, and one factor specific to each sector (land in farming,
capital in manufacturing), which cannot move. Each sector has constant returns
to scale and diminishing returns to labour; markets are competitive.
\end{definition}

Labour is allocated so that it earns the same in both sectors: $w = p_1\,
\mathrm{MPL}_1(L_1) = p_2\, \mathrm{MPL}_2(L - L_1)$. The diagram draws the two
value-of-marginal-product curves from the two ends of the labour force.

\begin{omfigure}
\begin{tikzpicture}[econ, x=0.75cm, y=0.6cm]
  \draw[econ axes] (0,0) -- (10,0);
  \draw[econ axes] (0,0) -- (0,9.3) node[above] {$w$};
  \draw[econ axes] (10,0) -- (10,9.3);
  \node[below] at (0,0) {$O_1$}; \node[below] at (10,0) {$O_2$};
  \draw[curve, omDef, domain=0:10, samples=40] plot (\x, {7*exp(-0.2*\x)}) node[right, font=\scriptsize] {$p_1\mathrm{MPL}_1$};
  \draw[curve, omDef, dashed, domain=0:10, samples=40] plot (\x, {9.1*exp(-0.2*\x)}) node[above right, font=\scriptsize] {$p_1'\mathrm{MPL}_1$};
  \draw[curve, omThm, domain=0:10, samples=40] plot (\x, {7*exp(-0.2*(10-\x))});
  \node[omThm, font=\scriptsize, anchor=south east] at (10,7.1) {$p_2\mathrm{MPL}_2$};
  \fill (5,2.575) circle (1.6pt); \fill (5.656,2.937) circle (1.6pt);
  \draw[guide] (5,0) node[below] {$L_1$} -- (5,2.575);
  \draw[guide] (5.656,0) node[below right] {$L_1'$} -- (5.656,2.937);
  \node[font=\scriptsize] at (2.5,-1.1) {labour in sector 1 $\rightarrow$};
  \node[font=\scriptsize] at (7.8,-1.1) {$\leftarrow$ labour in sector 2};
\end{tikzpicture}
\par\smallskip{\small The specific-factors diagram. A rise in the price of good 1
shifts its curve up in proportion; labour moves into sector 1 and the wage rises,
but by less than the price.}
\end{omfigure}

\begin{proposition}[A price rise in the specific-factors model]\label{prop:b3:factor-endowments:specific}
Let $p_1$ rise with $p_2$ fixed. Then $0 < \hat w < \hat p_1$; the real return to the
factor specific to sector 1 rises in terms of both goods, the real return to the
factor specific to sector 2 falls in terms of both, and the real wage rises in
terms of good 2 and falls in terms of good 1.
\end{proposition}

\begin{proof}
At the old allocation, $p_1\mathrm{MPL}_1$ exceeds the wage paid in sector 2, so
labour moves into sector 1 until the two are equal again: $L_1$ rises. Then
$\mathrm{MPL}_2$ rises (sector 2 has less labour), so $w = p_2\mathrm{MPL}_2$ rises:
$\hat w > 0$. And $\mathrm{MPL}_1$ falls, so $w = p_1\mathrm{MPL}_1$ rises by less
than $p_1$: $\hat w < \hat p_1$. The real wage $w/p_2$ rises and $w/p_1$ falls. With
constant returns, the return to the specific factor in sector $j$, in terms of
good $j$, is its marginal product, which rises with the labour working with it:
in sector 1 it rises, so the nominal return rises by more than $p_1$, and so more
than $p_2$; in sector 2 it falls, and with $p_2$ fixed and $p_1$ higher, it falls in
terms of both goods.
\end{proof}

Opening to trade raises the price of the export good. The model therefore
predicts a conflict by sector: owners of land and capital in the export
sector gain, those in the import-competing sector lose, and workers, who can
move, are caught in between.

\begin{model}{Specific factors}
\assumes In the short run some factors cannot move between sectors; labour can.
\predicts A trade shock divides a country by sector, not by factor: the owners
of the factors specific to an industry stand together for or against trade; the
wage moves by less than the price, so whether workers gain is ambiguous.
\tested In congressional testimony on a United States trade bill of 1973,
industries' business associations and unions took the same side on trade policy
in 19 of 21 industries. % ledger: b3fe:magee.lobby
The owners' side of the prediction fits; that the workers sided with their
industry fits a version of the model in which skills, too, are specific to an
industry.
\breaks In the long run capital is rebuilt and workers retrain, factors move
between sectors, and the division becomes one by factor: the subject of the
rest of the chapter.
\end{model}

\section{The Heckscher--Ohlin model}\label{sec:b3:factor-endowments:ho}

\begin{definition}[Factor intensity; factor abundance]\label{def:b3:factor-endowments:intensity}
\index{factor intensity}\index{factor abundance}
Good 1 is \emph{capital-intensive} relative to good 2 if at any given factor
prices it uses more capital per worker: $a_{K1}/a_{L1} > a_{K2}/a_{L2}$. A country is
\emph{capital-abundant} relative to another if it has more capital per worker:
$K/L > K^*/L^*$.
\end{definition}

\begin{definition}[Heckscher--Ohlin model]\label{def:b3:factor-endowments:ho}
\index{Heckscher--Ohlin model}
The \emph{Heckscher--Ohlin model} has two countries, two goods and two factors
that move freely between sectors but not between countries. Technologies are the
same in both countries, with constant returns to scale and a ranking of factor
intensities that does not reverse; tastes are the same and homothetic; markets
are competitive and trade is free.
\end{definition}

With identical technologies and tastes, the only difference between the two
countries is in their endowments. In autarky the capital-abundant country has
cheap capital, so the capital-intensive good is relatively cheap there; with
trade, it exports that good. This is the Heckscher--Ohlin prediction: \emph{a
country exports the goods that use its abundant factor intensively.}

\section{Factor prices}\label{sec:b3:factor-endowments:prices}

When both goods are made, competition drives profits to zero:
\begin{equation}\label{eq:b3:fe:zp}
  p_1 = a_{L1} w + a_{K1} r, \qquad p_2 = a_{L2} w + a_{K2} r .
\end{equation}
With fixed input requirements these are two linear equations in $w$ and $r$, and
goods prices alone determine factor prices. With cost-minimising requirements
that vary with factor prices, the same is true, and differentiating gives, by
the envelope theorem (the changes in the $a$'s have no first-order effect on
cost),
\begin{equation}\label{eq:b3:fe:hat}
  \hat p_j = \theta_{Lj}\hat w + \theta_{Kj}\hat r, \qquad j = 1, 2 .
\end{equation}

\begin{definition}[Magnification effect]\label{def:b3:factor-endowments:magnification}
\index{magnification effect}
A change in goods prices has a \emph{magnification effect} on factor prices if
one factor price rises proportionally more than both goods prices and the other
rises proportionally less than both (or falls).
\end{definition}

\begin{proposition}[Stolper--Samuelson theorem]\label{prop:b3:factor-endowments:ss}
If good 1 is labour-intensive and $\hat p_1 > \hat p_2$, then $\hat w > \hat p_1 > \hat p_2 >
\hat r$: the real wage rises in terms of both goods and the real rental falls in
terms of both.
\end{proposition}

\begin{proof}
By \eqref{eq:b3:fe:hat}, each $\hat p_j$ is a weighted average of $\hat w$ and $\hat r$,
with weights $\theta_{Lj}$ and $\theta_{Kj} = 1 - \theta_{Lj}$, so it lies between them.
Subtracting the two equations, $\hat p_1 - \hat p_2 = (\theta_{L1} - \theta_{L2})(\hat w -
\hat r)$. Good 1 is labour-intensive, so $\theta_{L1} > \theta_{L2}$, and $\hat p_1 > \hat p_2$
gives $\hat w > \hat r$. Both prices lie between $\hat r$ and $\hat w$, and $\hat p_1$, which
puts more weight on $\hat w$, lies closer to it: $\hat w > \hat p_1 > \hat p_2 > \hat r$.
\end{proof}

The theorem is the long-run counterpart of \cref{prop:b3:factor-endowments:specific}.
When trade raises the relative price of a country's exports, the factor used
intensively in them gains in real terms and the other factor loses, whatever
the sector its owners work in. In a capital-abundant country, opening to trade
lowers the real wage of its workers; in a labour-abundant one, it raises it.

\begin{omfigure}
\begin{tikzpicture}[econ, x=0.9cm, y=0.9cm]
  \draw[econ axes] (0,0) -- (8,0) node[right] {$w$};
  \draw[econ axes] (0,0) -- (0,5) node[above] {$r$};
  \draw[curve, omDef] (0,4) -- (5,0) node[below, font=\scriptsize] {$p_1$};
  \draw[curve, omDef, dashed] (0,4.4) -- (5.5,0) node[below right, font=\scriptsize] {$p_1'$};
  \draw[curve, omThm] (0,3) -- (7.5,0) node[below, font=\scriptsize] {$p_2$};
  \fill (2.5,2) circle (1.6pt); \fill (3.5,1.6) circle (1.6pt);
  \draw[guide] (2.5,0) -- (2.5,2) -- (0,2);
  \draw[guide] (3.5,0) -- (3.5,1.6) -- (0,1.6);
\end{tikzpicture}
\par\smallskip{\small Zero-profit lines $p_j = a_{Lj}w + a_{Kj}r$ for the two goods.
Their crossing gives the factor prices. Good 1, labour-intensive, has the
steeper line. A 10~\% rise in $p_1$ shifts it out (dashed): here the wage rises by
40~\% and the rental falls by 20~\%.}
\end{omfigure}

\begin{definition}[Factor-price equalisation]\label{def:b3:factor-endowments:fpe}
\index{factor-price equalisation}
\emph{Factor-price equalisation} is the prediction that free trade in goods
equalises factor prices across countries: if two countries face the same goods
prices, share the same technology and both make both goods, then
\eqref{eq:b3:fe:zp} gives them the same wage and the same rental.
\end{definition}

\omterm{def:b3:factor-endowments:fpe}{Factor-price equalisation} is the model's most striking prediction and its most
obviously false one: wages for the same work differ many times over between
countries. The gap is the measure of how much the model leaves out ---
technology above all (\cref{exo:b3:factor-endowments:9}).

The Stolper--Samuelson prediction (\cref{prop:b3:factor-endowments:ss}) is easier to test inside one country: when
the price of goods made with unskilled labour falls, unskilled wages should fall
relative to others. Whether imports from low-wage countries did that to unskilled
American wages is the trade-and-wages debate of \cref{sec:b3:trade-policy:evidence},
where the evidence is weighed; the price changes that the theorem needs turned out
to be small, and the effects most studies find run through local labour markets,
which the model of this chapter does not have.

\section{Outputs: Rybczynski}\label{sec:b3:factor-endowments:rybczynski}

\begin{definition}[Rybczynski effect]\label{def:b3:factor-endowments:rybczynski}
\index{Rybczynski effect}
The \emph{Rybczynski effect} is the response of outputs to a change in a
country's endowments at given goods prices.
\end{definition}

\begin{proposition}[Rybczynski theorem]\label{prop:b3:factor-endowments:rybczynski}
At given goods prices, with both goods made, a rise in the endowment of capital
raises the output of the capital-intensive good more than in proportion and
lowers the output of the other good.
\end{proposition}

\begin{proof}
Fixed goods prices fix factor prices by \eqref{eq:b3:fe:zp}, and so the input
requirements. Full employment then reads $a_{K1}Q_1 + a_{K2}Q_2 = K$ and $a_{L1}Q_1 +
a_{L2}Q_2 = L$. Let good 2 be capital-intensive and $\Delta = a_{L1}a_{K2} -
a_{K1}a_{L2} > 0$. By Cramer's rule, $Q_2 = (a_{L1}K - a_{K1}L)/\Delta$ and $Q_1 = (a_{K2}L -
a_{L2}K)/\Delta$. With $L$ fixed, $\partial Q_1/\partial K = -a_{L2}/\Delta < 0$. And
$\hat Q_2 = a_{L1}K\hat K/(a_{L1}K - a_{K1}L) > \hat K$, since the denominator is smaller
than the numerator's $a_{L1}K$.
\end{proof}

An economy that accumulates capital, at given world prices, does not spread it
over all sectors: it moves workers out of labour-intensive sectors to staff the
capital-intensive ones, and factor prices do not move at all. Immigration works
the same way with the roles reversed (\cref{exo:b3:factor-endowments:12}).

\section{The factor content of trade and its tests}\label{sec:b3:factor-endowments:content}

With many goods and many factors, the Heckscher--Ohlin prediction about goods
becomes a prediction about the factors that goods embody.

\begin{definition}[Factor content of trade; Heckscher--Ohlin--Vanek equation; missing trade]\label{def:b3:factor-endowments:factor-content}
\index{factor content of trade}\index{Heckscher--Ohlin--Vanek equation}\index{missing trade}
The \emph{factor content of trade} of country $c$ is $F_c = A T_c$, the amounts
of each factor used to make its net exports, where column $j$ of $A$ gives the
factor requirements of good $j$. The \emph{Heckscher--Ohlin--Vanek (HOV)
equation} predicts $F_c = V_c - s_c V_w$, where $V_w$ is the world's endowments.
\emph{Missing trade} is the finding that measured factor contents are much
smaller than the HOV prediction.
\end{definition}

\begin{proposition}[The HOV equation]\label{prop:b3:factor-endowments:hov}
If all countries share the technology $A$, factor prices are equal, every
country makes goods with full employment ($A Q_c = V_c$), and tastes are identical
and homothetic, so that every country consumes the same proportions of the
world's output ($C_c = s_c Q_w$), then $F_c = V_c - s_c V_w$.
\end{proposition}

\begin{proof}
$F_c = A T_c = A(Q_c - C_c) = AQ_c - s_c AQ_w$. Full employment in country $c$ gives
$AQ_c = V_c$, and summing over countries, $AQ_w = V_w$. Hence $F_c = V_c - s_c V_w$.
\end{proof}

A country is a net exporter of the services of a factor exactly when its share
of the world's endowment of that factor exceeds its share of world spending.
Read the opening chart with output standing in for spending. Japan, France and
Italy have larger shares of the world's capital than of its output, and should
export capital services embodied in goods; India and China, the services of
labour. But the United States has smaller shares of both the world's capital and
its workers than of its output, and China larger shares of both: measured this
way, the richest large economy is scarce in every factor and a poorer one abundant
in all. That pattern is the first clue in the record below.

\begin{record}[Leontief's paradox]
The first test came before the HOV equation was written. Leontief used the
input--output table of the United States for 1947 to compute the capital and labour
needed to make a million dollars' worth of its exports, and a million dollars' worth
of the goods that compete with its imports. The trade of the most capital-rich
economy of the time turned out to be, on balance, labour-using and
capital-conserving. Leontief's own explanation was that an American worker was as
efficient as three foreign ones. % ledger: b3fe:ken.leontief
In a second study he sorted the workers by occupation. Export production used more
skill: % ledger: b3fe:ken.occ
\begin{center}
\small
\begin{tabular}{lcc}
  Percent of workers & Export production & Import-competing production \\ \hline
  Professional and technical & 13.75 & 12.24 \\
  Clerical, sales and service & 22.07 & 17.00 \\
  Craftsmen and foremen & 15.15 & 11.79 \\
  Operatives & 30.05 & 28.38 \\
  Labourers & 18.98 & 30.59 \\ \hline
\end{tabular}
\end{center}
Kenen counted skill as capital: he valued each worker's wage above a labourer's as
the return on human capital and added it to the physical capital. Total capital per
worker in exports and in import-competing production then came to 30.61 and 29.83
thousand 1947 dollars with the wage differences capitalised at 9~\%, 27.16 and 26.89 at 11~\%, and
25.08 and 25.12 at 12.7~\%: the paradox reverses in two of the three cases. Physical capital alone
still runs the paradoxical way. % ledger: b3fe:ken.total
Leamer (1980) replied that Leontief had made the wrong comparison: when trade is not
balanced, exports set against imports need not reveal abundance. % ledger: b3fe:leamer
\end{record}

The proposition below makes Leamer's point exact.

\begin{proposition}[Leamer's revealed abundance]\label{prop:b3:factor-endowments:leamer}
Under the assumptions of \cref{prop:b3:factor-endowments:hov}, a country is abundant in
capital relative to labour ($K_c/L_c > K_w/L_w$) if and only if the capital per worker
embodied in its production exceeds the capital per worker embodied in its
consumption. The capital per worker embodied in its exports and in its imports can
rank the other way.
\end{proposition}

\begin{proof}
Production embodies $V_c = (K_c, L_c)$. Consumption embodies $V_c - F_c = s_c V_w$, whose
capital per worker is $K_w/L_w$. So production is more capital-intensive than
consumption exactly when $K_c/L_c > K_w/L_w$. For the second claim, write $F = A X - A M$ for
exports $X$ and imports $M$. When a trade surplus makes both entries of $F$ positive,
$(A X)_K / (A X)_L$ is a weighted average of $(A M)_K/(A M)_L$ and $F_K/F_L$, so it is
below the imports' ratio whenever $F_K / F_L$ is; the listing below gives numbers.
\end{proof}

\omcode{code/bachelor-3/08-factor-endowments/listings/leontief_leamer.py}{1}{21}{Leontief's comparison and Leamer's on a country with a trade surplus.}\label{lst:b3:factor-endowments:leamer}

\omoutput{code/bachelor-3/08-factor-endowments/outputs/leontief_leamer.txt}{Factor units per year.}\label{out:b3:factor-endowments:leamer}

Home's exports carry 3.25 units of capital per worker, its imports 4.0: by Leontief's
test it is abundant in labour. Yet its production carries 2.16 against 2.0 in its
consumption and in the world: it is abundant in capital. % output: leontief_leamer
It exports labour services on net only because it spends less than it earns: its
consumption share is 0.2 of the world's, below its share of the world's labour (0.25),
so it sells the services of both factors abroad. A test of the
model must compare what a country makes with what it consumes, which is what the
HOV equation does.

\begin{record}[The missing trade]
Bowen, Leamer and Sveikauskas tested the HOV equations with twelve resources and
the trade of twenty-seven countries in 1967; the data did not support the
prediction that trade reveals \omterm{def:b3:factor-endowments:intensity}{factor abundance}. % ledger: b3fe:bls.tests
Trefler found that trade is \emph{missing}: measured factor contents are far
smaller than the predicted ones, and rich countries appear scarce in most
factors while poor ones appear abundant in all. He traced the failures to home
bias in consumption and to differences in technology. % ledger: b3fe:tre.missing
Davis and Weinstein then allowed technology to differ across countries, as
measured in each country's own input data, together with trade costs and
\omterm{def:b3:ricardian-trade:iceberg}{non-traded goods}: with these amendments, countries export the services of their
abundant factors in about the right amounts. % ledger: b3fe:dw.works
\end{record}

The repairs bring Ricardo back. Technology does differ: the chart below shows
how weakly capital per worker alone accounts for output per worker.

\begin{omchart}{Output per worker against capital per worker, 2019, 175
countries (thousands of 2021 dollars per person engaged, log scales), with the
least-squares line in logs.}{Penn World Table 11.0; our computation}
\begin{tikzpicture}
\begin{axis}[omscatter, width=0.7\linewidth, height=0.45\linewidth, xmode=log, ymode=log,
  xmin=1, xmax=2000, ymin=1, ymax=300, log ticks with fixed point,
  xlabel={Capital per worker}, ylabel={Output per worker}]
  \addplot table[x=k, y=y] {figdata/bachelor-3/08-factor-endowments/out/kpw.csv};
  \addplot[omThm, thick, no marks, sharp plot] table[x=k, y=y] {figdata/bachelor-3/08-factor-endowments/out/kpw_fit.csv};
\end{axis}
\end{tikzpicture}
\end{omchart}

Capital per worker is about 30 times as high at the 90th percentile of
countries as at the 10th, % figdata: kpw_stats
and output per worker rises with it with an elasticity of 0.76 --- above the
median share of capital in income in the same table, 0.48. If technology were the
same everywhere, the elasticity would be that share; it is larger because
countries with more capital also have better technology and more skills. The \omterm{def:b3:factor-endowments:ho}{Heckscher--Ohlin model}, which
assumes the same technology everywhere, holds only once those differences are
measured.

\begin{model}{Heckscher--Ohlin}
\assumes Identical technologies and homothetic tastes everywhere; two or more
factors, mobile between sectors, immobile between countries; free trade.
\predicts Countries export goods intensive in their abundant factors; trade
raises the real return of the abundant factor and lowers that of the scarce one
(Stolper--Samuelson); factor prices equalise; the \omterm{def:b3:factor-endowments:factor-content}{factor content of trade}
equals $V_c - s_c V_w$.
\tested In its pure form HOV is rejected: the data do not support the prediction
that trade reveals \omterm{def:b3:factor-endowments:intensity}{factor abundance}, and trade is missing. With measured
technology differences, home bias and trade costs, the \omterm{def:b3:factor-endowments:factor-content}{factor content of trade}
lines up with endowments in sign and size.
\breaks Technologies differ across countries, so factor prices do not equalise;
consumers favour home goods; and trade costs keep many goods at home.
\end{model}

\section{Exercises}\label{sec:b3:factor-endowments:exercises}

\begin{exercise}[$\star$]\label{exo:b3:factor-endowments:1}
In the opening chart, which countries have a larger share of world capital
than of world output? By HOV, which factor's services should they export?
\end{exercise}

\begin{exercise}[$\star$]\label{exo:b3:factor-endowments:2}
Making a million dollars of steel uses \$4 million of capital and 20 workers;
a million dollars of clothing uses \$1 million of capital and 40 workers. Which
good is capital-intensive? Is that ranking still true if wages double?
\end{exercise}

\begin{exercise}[$\star$]\label{exo:b3:factor-endowments:3}
Good 1 needs 2 workers and 1 unit of capital per unit; good 2 needs 1 worker and 2
units of capital. If $p_1 = 10$ and $p_2 = 8$, find $w$ and $r$.
\end{exercise}

\begin{exercise}[$\star$]\label{exo:b3:factor-endowments:4}
A labour-abundant country's exports are found to be more capital-intensive than
its imports. Does this contradict the \omterm{def:b3:factor-endowments:ho}{Heckscher--Ohlin model}? What else would you
check?
\end{exercise}

\begin{exercise}[$\star\star$]\label{exo:b3:factor-endowments:5}
Labour's cost share is 0.7 in good 1 and 0.4 in good 2. The price of good 1 rises
by 10~\%, that of good 2 is unchanged. Find $\hat w$ and $\hat r$, and check the
\omterm{def:b3:factor-endowments:magnification}{magnification effect}.
\end{exercise}

\begin{exercise}[$\star\star$]\label{exo:b3:factor-endowments:6}
With the technology of \cref{exo:b3:factor-endowments:3}, a country has $L = 100$
and $K = 80$. Find its outputs. Capital rises by 10~\%: find the new outputs and
their percentage changes.
\end{exercise}

\begin{exercise}[$\star\star$]\label{exo:b3:factor-endowments:7}
In Brenland, grain is made with land and labour, manufactures with capital and
labour, and workers move freely. The world price of grain rises by 20~\%. Who gains
and who loses in the short run? In the long run, if capital can be turned into
farm equipment?
\end{exercise}

\begin{exercise}[$\star\star$]\label{exo:b3:factor-endowments:8}
Three goods (steel, cars, clothing) need 4, 2 and 1 units of capital and 1, 2 and
3 units of labour per unit. A country's net exports are 10, $-5$ and $-8$ units, and
its endowments are 400 units of capital and 500 of labour. Compute the factor
content of its trade and compare it with its endowments.
\end{exercise}

\begin{exercise}[$\star\star$]\label{exo:b3:factor-endowments:9}
What wage data would contradict \omterm{def:b3:factor-endowments:fpe}{factor-price equalisation}? Do they reject the
\omterm{def:b3:factor-endowments:ho}{Heckscher--Ohlin model}, or one of its assumptions? Which?
\end{exercise}

\begin{exercise}[$\star\star\star$]\label{exo:b3:factor-endowments:10}
Prove \omterm{def:b3:factor-endowments:fpe}{factor-price equalisation} with fixed input requirements: show that if
$a_{L1}a_{K2} \ne a_{K1}a_{L2}$, the zero-profit conditions have exactly one solution
$(w, r)$, and state what is needed for it to be positive.
\end{exercise}

\begin{exercise}[$\star\star\star$]\label{exo:b3:factor-endowments:11}
Simulate an HOV world of ten countries and three factors. Compare measured
factor contents with the HOV prediction when (a) HOV holds exactly; (b) four
fifths of the predicted trade is missing; (c) trade is missing and measured with
noise. Report the share of matching signs and the slope of measured on
predicted contents. What does each test detect?
\end{exercise}

\begin{exercise}[$\star\star\star$]\label{exo:b3:factor-endowments:12}
In a small open Heckscher--Ohlin economy that makes both goods, immigration raises
the labour force by 5~\%. Show that wages need not fall, and state the conditions:
on prices, on the economy's size, and on how large the inflow can be.
\end{exercise}

\section{Problem: Kestria's Capital Boom}\label{sec:b3:factor-endowments:problem}

\begin{problem}[{Weekend problem --- Kestria's capital boom}]\label{pb:b3:factor-endowments:boom}
Kestria and Brenland (invented) make machines and textiles with capital and
labour, with the same fixed input requirements per unit:
% figdata: boom_params
\begin{center}\small
\pgfplotstabletypeset[col sep=comma,
  columns/item/.style={column name={Good}, string type, column type=l},
  columns/capital/.style={column name={Capital per unit}, fixed, precision=0},
  columns/labour/.style={column name={Workers per unit}, fixed, precision=0},
  columns/price/.style={column name={World price (kests)}, fixed, precision=0},
  every head row/.style={after row=\hline}, every last row/.style={after row=\hline}]
  {figdata/bachelor-3/08-factor-endowments/out/boom_params.csv}
\end{center}
Kestria has 60 million units of capital and 40 million workers; Brenland 30
million and 50 million; the world as a whole 1{,}300 million and 1{,}100 million.
Consumers everywhere spend half of their income on each good. Both countries are
small, and with trade they face the world prices in the table.

\medskip\noindent\textbf{Part I --- Kestria in autarky.}
\begin{enumerate}
\item Which good is capital-intensive?
\item Compare capital per worker in Kestria, Brenland and the world. Which country
is capital-abundant?
\item Find Kestria's outputs with all of its capital and labour employed.
\item In autarky, equal spending on the two goods sets the relative price of
machines equal to $Q_T/Q_M$. Compute it.
\item With the price of textiles set to 1, find Kestria's autarky wage and
rental.
\item Check that Kestria's output, valued at \omterm{def:b3:ricardian-trade:autarky}{autarky prices}, equals its factor
income.
\end{enumerate}

\medskip\noindent\textbf{Part II --- free trade.}
\begin{enumerate}[resume]
\item Find the wage and the rental at world prices.
\item Find Brenland's outputs. Why are its factor prices the same as Kestria's?
\item Find Kestria's consumption and net exports of each good, and check that
trade balances.
\item Compute Kestria's real wage and real rental, in machines and in textiles,
in autarky and with trade.
\item By what percentage does the relative price of machines rise, and the
wage and the rental in terms of textiles? Check the \omterm{def:b3:factor-endowments:magnification}{magnification effect}.
\item Who in Kestria gains from trade and who loses?
\end{enumerate}

\medskip\noindent\textbf{Part III --- the capital boom.}
\begin{enumerate}[resume]
\item Kestria's capital stock rises by 25~\%, to 75 million. At world prices,
find its new outputs.
\item Give the percentage changes of the two outputs, and relate them to the
Rybczynski theorem.
\item Why does textile output fall when no price has changed?
\item What happens to the wage and the rental? To workers' real wage?
\item Find Kestria's new GDP, consumption and net exports.
\item Why is it reasonable to hold world prices fixed? What would change if
Kestria were large?
\end{enumerate}

\medskip\noindent\textbf{Part IV --- the factor content of trade.}
\begin{enumerate}[resume]
\item Before the boom, compute the factor content of Kestria's net exports.
\item Compute Kestria's share of world spending.
\item Check the HOV equation for Kestria.
\item Do the same for Brenland.
\item What would ``\omterm{def:b3:factor-endowments:factor-content}{missing trade}'' look like in these data, and what could cause
it here?
\item State the results: the change in Kestria's textile output and in its
workers' real wage after the boom.
\end{enumerate}
\end{problem}
